\section{Роли и ответственность}

\begin{frame}{Обеспечивает / не делает}
  \vspace{0.3em}
  \begin{columns}[T]
    \begin{column}{0.49\textwidth}
      \begin{block}{Обеспечивает}
        \begin{itemize}\setlength\itemsep{0.28em}
          \item пункт ответственности;
          \item второй пункт;
          \item \textbf{важное условие} выделено жирным;
          \item заключительный пункт списка.
        \end{itemize}
      \end{block}
    \end{column}
    \begin{column}{0.49\textwidth}
      \begin{block}{Не делает}
        \begin{itemize}\setlength\itemsep{0.28em}
          \item не входит в зону ответственности;
          \item не принимает решение X;
          \item не выполняет функцию Y.
        \end{itemize}
      \end{block}
      \begin{block}{Центральный орган}
        Кратко: что делает вышестоящий или координирующий орган.
      \end{block}
    \end{column}
  \end{columns}
\end{frame}
